\documentclass[10pt,a4paper,twocolumn]{article}

% ==================================================
% Packages
% ==================================================

\usepackage{xepersian}
\usepackage{fontspec}
\usepackage{amsmath}
\usepackage{amssymb}
\usepackage{graphicx}
\usepackage{geometry}
\usepackage{fancyhdr}
\usepackage{titlesec}
\usepackage{hyperref}
\usepackage{enumitem}
\usepackage{float}
\usepackage{adjustbox}

% ==================================================
% Page settings
% ==================================================

\geometry{
    a4paper,
    top=2cm,
    bottom=2cm,
    left=1.6cm,
    right=1.6cm,
    columnsep=0.7cm
}

% ==================================================
% Fonts
% ==================================================

\settextfont{Vazirmatn}
\setdigitfont{Vazirmatn}
\setlatintextfont{Times New Roman}

% ==================================================
% Header and Footer
% ==================================================

\pagestyle{fancy}
\fancyhf{}
\fancyhead[L]{گزارش پروژه}
\fancyhead[R]{یادگیری ماشین و هوش مصنوعی}
\fancyfoot[C]{\thepage}

\setlength{\headheight}{13pt}
\renewcommand{\headrulewidth}{0.4pt}
\renewcommand{\footrulewidth}{0.4pt}

% ==================================================
% Section formatting
% ==================================================

\titleformat{\section}
{\normalfont\bfseries}
{\thesection.}
{0.5em}
{}

\titleformat{\subsection}
{\normalfont\bfseries}
{\thesubsection.}
{0.5em}
{}

\titlespacing{\section}
{0pt}
{10pt}
{5pt}

\titlespacing{\subsection}
{0pt}
{7pt}
{3pt}

% ==================================================
% Paragraph settings
% ==================================================

\setlength{\parindent}{1em}
\setlength{\parskip}{0pt}

% ==================================================
% Hyperlinks
% ==================================================

\hypersetup{
    colorlinks=true,
    linkcolor=black,
    citecolor=black,
    urlcolor=black
}

% ==================================================
% Document
% ==================================================

\begin{document}

% ==================================================
% Title
% ==================================================

\begin{center}

\vspace{0.7cm}

{\Large\textbf{
بررسی روش‌های یادگیری ماشین در تشخیص تقلب مالی
}}

\vspace{0.5cm}

{\small
نام نویسنده: امیرحسین رنجبر
}

\vspace{0.15cm}

\end{center}

\vspace{0.3cm}

% ==================================================
% Abstract
% ==================================================

\section*{چکیده}

تشخیص تراکنش‌های مالی تقلبی یکی از مسائل مهم در حوزه
یادگیری ماشین است. نامتوازن بودن داده‌ها، تعداد بسیار کم
نمونه‌های تقلبی و هزینه بالای تشخیص نادرست از جمله چالش‌های
اصلی این مسئله هستند.

در این گزارش، ابتدا چند مدل پایه شامل درخت تصمیم،
KNN و رگرسیون لجستیک بررسی می‌شوند. سپس مدل‌های پیشرفته‌تر
شامل XGBoost و جنگل تصادفی مورد بررسی قرار می‌گیرند.
همچنین مدل‌های مختلف در یک روش Ensemble ترکیب می‌شوند.

عملکرد مدل‌ها با استفاده از معیارهایی مانند Precision،
Recall، \lr{F1-score} و PR-AUC ارزیابی می‌شود. همچنین برای
انتخاب تنظیمات مناسب مدل‌ها از روش‌های اعتبارسنجی متقابل
و بهینه‌سازی هایپرپارامترها استفاده می‌شود.

\textbf{واژگان کلیدی:}

یادگیری ماشین، تشخیص تقلب، تراکنش مالی، درخت تصمیم،
KNN، رگرسیون لجستیک، جنگل تصادفی، XGBoost، Ensemble،
داده‌های نامتوازن

% ==================================================
% 1. Introduction
% ==================================================

\section{مقدمه}

افزایش استفاده از سامانه‌های پرداخت الکترونیکی باعث شده است
که حجم بسیار زیادی از تراکنش‌های مالی در هر روز ثبت شود.
در کنار این افزایش، احتمال وقوع تراکنش‌های تقلبی نیز وجود دارد.
تشخیص خودکار تراکنش‌های تقلبی می‌تواند نقش مهمی در کاهش
خسارت‌های مالی داشته باشد. با این حال، مسئله تشخیص تقلب
معمولاً با مشکل نامتوازن بودن کلاس‌ها مواجه است.

در بسیاری از مجموعه‌داده‌های واقعی، تعداد تراکنش‌های سالم
بسیار بیشتر از تراکنش‌های تقلبی است. بنابراین یک مدل ممکن است
دقت بسیار بالایی داشته باشد، اما در تشخیص نمونه‌های تقلبی
عملکرد مناسبی نداشته باشد.

به همین دلیل، در این مسئله معیارهایی مانند Precision،
Recall، \lr{F1-score} و PR-AUC اهمیت زیادی پیدا می‌کنند.

در این پژوهش، ابتدا مدل‌های پایه بررسی شده و سپس مدل‌های
پیشرفته‌تر و در نهایت روش Ensemble مورد استفاده قرار گرفته‌اند.

% ==================================================
% 2. Dataset
% ==================================================

\section{معرفی داده‌ها}

در این پژوهش از یک مجموعه‌داده شامل تراکنش‌های مالی استفاده
شده است. هر نمونه شامل تعدادی ویژگی مربوط به تراکنش و یک
برچسب دودویی است.

برچسب صفر نشان‌دهنده تراکنش سالم و برچسب یک نشان‌دهنده
تراکنش تقلبی است.

یکی از ویژگی‌های مهم داده، عدم توازن شدید بین دو کلاس است.
به همین دلیل، استفاده از Accuracy به‌تنهایی نمی‌تواند معیار
مناسبی برای ارزیابی مدل باشد.

داده‌ها پس از انجام مراحل پیش‌پردازش به دو مجموعه آموزش و
آزمون تقسیم شدند.

% ==================================================
% 3. Data preprocessing
% ==================================================

\section{پیش‌پردازش داده‌ها}

پیش از آموزش مدل‌ها، داده‌ها بررسی و آماده‌سازی شدند.

مراحل پیش‌پردازش شامل بررسی مقادیر گمشده، داده‌های تکراری
و آماده‌سازی ویژگی‌ها برای ورود به مدل است.

همان‌طور که در بررسی اولیه داده‌ها مشاهده شد، مجموعه‌داده
دارای تعدادی مقدار تکراری بود. از آنجا که وجود داده‌های
تکراری می‌تواند بر فرآیند آموزش مدل تأثیر بگذارد، این
نمونه‌ها شناسایی و حذف شدند.

تعداد نمونه‌ها و تراکنش‌های تقلبی قبل و بعد از حذف داده‌های
تکراری در جدول زیر ارائه شده است.

\begin{table}[H]

\centering

\caption{مقایسه داده‌ها قبل و بعد از حذف نمونه‌های تکراری}

\vspace{0.15cm}

\begin{adjustbox}{max width=\columnwidth}

\begin{tabular}{|l|r|r|}

\hline

\textbf{Data} & \textbf{Before} & \textbf{After} \\

\hline

Total transactions & 284807 & 284315 \\

\hline

Fraudulent transactions & 492 & 473 \\

\hline

\end{tabular}

\end{adjustbox}

\label{tab:duplicate-removal}

\end{table}

همان‌طور که مشاهده می‌شود، تعداد کل تراکنش‌ها از
284807 به 284315 کاهش یافته است. همچنین تعداد تراکنش‌های
تقلبی از 492 نمونه به 473 نمونه رسیده است.

داده‌ها به دو مجموعه آموزش و آزمون تقسیم شدند:

\[
D = D_{\mathrm{train}} \cup D_{\mathrm{test}}
\]

به‌طوری‌که مجموعه آزمون در فرآیند آموزش مدل استفاده نمی‌شود.

در صورت وجود ویژگی‌هایی با مقیاس‌های متفاوت، می‌توان از
روش‌های استانداردسازی برای آماده‌سازی داده‌ها استفاده کرد.
استانداردسازی به‌خصوص برای مدل‌هایی مانند KNN و رگرسیون
لجستیک اهمیت دارد.

% ==================================================
% 4. Base Models
% ==================================================

\section{مدل‌های پایه}

در مرحله اول پژوهش، سه مدل پایه برای ایجاد یک معیار اولیه
از عملکرد سیستم مورد بررسی قرار گرفتند:

\begin{itemize}
    \item درخت تصمیم
    \item KNN
    \item رگرسیون لجستیک
\end{itemize}

هدف از استفاده از مدل‌های پایه، بررسی عملکرد الگوریتم‌های
مختلف و ایجاد مبنایی برای مقایسه با مدل‌های پیشرفته‌تر است.

% ==================================================
% Decision Tree
% ==================================================

\subsection{درخت تصمیم}

درخت تصمیم یکی از الگوریتم‌های یادگیری ماشین است که با تقسیم
متوالی داده‌ها بر اساس ویژگی‌های مختلف، ساختاری درختی برای
انجام پیش‌بینی ایجاد می‌کند.

پیچیدگی و عملکرد این مدل به هایپرپارامترهای مختلفی وابسته است.
در این آزمایش، سه هایپرپارامتر مهم شامل
\lr{max\_depth}،
\lr{min\_samples\_split} و
\lr{min\_samples\_leaf} مورد بررسی قرار گرفتند.

هایپرپارامتر \lr{max\_depth} حداکثر عمق درخت را مشخص می‌کند.
افزایش این مقدار امکان ایجاد ساختارهای پیچیده‌تر را فراهم
می‌کند، اما افزایش بیش از حد آن می‌تواند باعث بیش‌برازش شود.

برای این پارامتر مقادیر
$3$، $5$، $8$، $10$، $15$، $20$ و \lr{None}
بررسی شدند.

پارامتر \lr{min\_samples\_split} حداقل تعداد نمونه‌های مورد
نیاز برای انجام تقسیم یک گره را مشخص می‌کند. مقادیر
$2$، $5$، $10$ و $20$ برای این پارامتر در نظر گرفته شدند.

پارامتر \lr{min\_samples\_leaf} حداقل تعداد نمونه‌هایی را
مشخص می‌کند که باید در هر گره برگ قرار داشته باشند. مقادیر
$1$، $2$، $5$ و $10$ برای این پارامتر بررسی شدند.

برای انتخاب بهترین ترکیب از این هایپرپارامترها، از روش
\lr{Grid Search} استفاده شد.

از آنجا که برای سه هایپرپارامتر به ترتیب $7$، $4$ و $4$ مقدار
در نظر گرفته شد، در مجموع:

\[
7 \times 4 \times 4 = 112
\]

ترکیب مختلف مورد بررسی قرار گرفت.

برای ارزیابی هر ترکیب از \lr{Stratified K-Fold Cross Validation}
با پنج بخش استفاده شد. معیار \lr{F1-score} به عنوان معیار
انتخاب بهترین ترکیب در نظر گرفته شد.

در نهایت، بهترین ترکیب هایپرپارامترها به صورت زیر به دست آمد:

\begin{table}[H]

\centering

\caption{هایپرپارامترهای انتخاب‌شده برای درخت تصمیم}

\vspace{0.15cm}

\begin{adjustbox}{max width=\columnwidth}

\begin{tabular}{|c|c|c|}

\hline

\textbf{هایپرپارامتر} &
\textbf{مقادیر بررسی‌شده} &
\textbf{مقدار انتخاب‌شده} \\

\hline

\lr{max\_depth} &
$3,5,8,10,15,20,None$ &
$8$ \\

\hline

\lr{min\_samples\_split} &
$2,5,10,20$ &
$20$ \\

\hline

\lr{min\_samples\_leaf} &
$1,2,5,10$ &
$5$ \\

\hline

\end{tabular}

\end{adjustbox}

\label{tab:dt_best_params}

\end{table}

پس از انتخاب هایپرپارامترها، مدل نهایی با استفاده از داده‌های
آموزشی آموزش داده شد و سپس روی مجموعه داده‌های آزمون ارزیابی
گردید.

\begin{table}[H]

\centering

\caption{نتایج مدل درخت تصمیم روی مجموعه آزمون}

\vspace{0.15cm}

\begin{adjustbox}{max width=\columnwidth}

\begin{tabular}{|c|c|c|c|c|}

\hline

\textbf{کلاس} &
\textbf{Precision} &
\textbf{Recall} &
\textbf{\lr{F1-score}} &
\textbf{Support} \\

\hline

$0$ & $1.00$ & $1.00$ & $1.00$ & $56651$ \\

\hline

$1$ & $0.92$ & $0.71$ & $0.80$ & $95$ \\

\hline

\textbf{Macro Avg} &
$0.96$ &
$0.85$ &
$0.90$ &
$56746$ \\

\hline

\textbf{Weighted Avg} &
$1.00$ &
$1.00$ &
$1.00$ &
$56746$ \\

\hline

\end{tabular}

\end{adjustbox}

\label{tab:dt_results}

\end{table}

همچنین ماتریس درهم‌ریختگی مدل به صورت زیر به دست آمد:

\begin{table}[H]

\centering

\caption{ماتریس درهم‌ریختگی مدل درخت تصمیم}

\vspace{0.15cm}

\begin{adjustbox}{max width=\columnwidth}

\begin{tabular}{|c|c|c|}

\hline

&
\textbf{پیش‌بینی کلاس 0} &
\textbf{پیش‌بینی کلاس 1} \\

\hline

\textbf{واقعی کلاس 0} &
$56645$ &
$6$ \\

\hline

\textbf{واقعی کلاس 1} &
$28$ &
$67$ \\

\hline

\end{tabular}

\end{adjustbox}

\label{tab:dt_confusion}

\end{table}

بر اساس نتایج، مدل برای کلاس اقلیت دارای
\lr{Precision} برابر با $0.92$ و
\lr{Recall} برابر با $0.71$ است.

مدل توانسته است $67$ مورد از مجموع $95$ نمونه مثبت را به درستی
شناسایی کند، در حالی که $28$ نمونه مثبت به اشتباه به عنوان
کلاس $0$ پیش‌بینی شده‌اند.

همچنین تنها $6$ نمونه از کلاس $0$ به اشتباه به عنوان کلاس $1$
شناسایی شده‌اند. مقدار \lr{F1-score} برابر با $0.80$ است.

% ==================================================
% KNN
% ==================================================

\subsection{K نزدیک‌ترین همسایه}

الگوریتم KNN یکی از روش‌های مبتنی بر فاصله است که برای
طبقه‌بندی یک نمونه جدید، نزدیک‌ترین نمونه‌های موجود در
داده‌های آموزشی را بررسی می‌کند.

پارامتر \lr{n\_neighbors} یا همان $K$، تعداد همسایه‌هایی را
مشخص می‌کند که در تصمیم‌گیری مورد استفاده قرار می‌گیرند.

از آنجا که KNN بر اساس فاصله بین نمونه‌ها عمل می‌کند،
استانداردسازی ویژگی‌ها اهمیت زیادی دارد. در این پژوهش،
داده‌های آموزشی پیش از استفاده در KNN استانداردسازی شدند.

برای انتخاب بهترین تنظیمات مدل، از \lr{Grid Search} به همراه
\lr{Stratified K-Fold Cross Validation} استفاده شد.

سه هایپرپارامتر شامل \lr{n\_neighbors}،
\lr{weights} و \lr{p} مورد بررسی قرار گرفتند.

برای \lr{n\_neighbors} مقادیر
$3$، $5$، $7$، $9$، $11$، $15$، $20$، $25$ و $30$
بررسی شدند.

در پارامتر \lr{weights} دو حالت \lr{uniform} و
\lr{distance} بررسی شد. در حالت \lr{distance}، نمونه‌های
نزدیک‌تر وزن بیشتری در تصمیم‌گیری دارند.

پارامتر \lr{p} نوع فاصله مورد استفاده را مشخص می‌کند.
مقدار $p=1$ متناظر با فاصله Manhattan و مقدار $p=2$
متناظر با فاصله Euclidean است.

در مجموع:

\[
9 \times 2 \times 2 = 36
\]

ترکیب مختلف بررسی شد.

برای ارزیابی هر ترکیب از \lr{Stratified K-Fold Cross Validation}
با پنج بخش استفاده شد و معیار انتخاب بهترین ترکیب
\lr{F1-score} بود.

\begin{table}[H]

\centering

\caption{هایپرپارامترهای انتخاب‌شده برای KNN}

\vspace{0.15cm}

\begin{adjustbox}{max width=\columnwidth}

\begin{tabular}{|c|c|c|}

\hline

\textbf{هایپرپارامتر} &
\textbf{مقادیر بررسی‌شده} &
\textbf{مقدار انتخاب‌شده} \\

\hline

\lr{n\_neighbors} &
$3,5,7,9,11,15,20,25,30$ &
$3$ \\

\hline

\lr{weights} &
\lr{uniform}, \lr{distance} &
\lr{distance} \\

\hline

\lr{p} &
$1,2$ &
$1$ \\

\hline

\end{tabular}

\end{adjustbox}

\label{tab:knn_best_params}

\end{table}

در نتیجه، بهترین تنظیمات KNN شامل $K=3$، وزن‌دهی
\lr{distance} و فاصله Manhattan ($p=1$) بود.

مقدار بهترین \lr{F1-score} در فرآیند اعتبارسنجی برابر با
$0.8543$ به دست آمد.

پس از انتخاب هایپرپارامترها، مدل نهایی با استفاده از داده‌های
آموزشی استانداردشده آموزش داده شد و روی مجموعه آزمون ارزیابی
گردید.

\begin{table}[H]

\centering

\caption{نتایج مدل KNN روی مجموعه آزمون}

\vspace{0.15cm}

\begin{adjustbox}{max width=\columnwidth}

\begin{tabular}{|c|c|c|c|c|}

\hline

\textbf{کلاس} &
\textbf{Precision} &
\textbf{Recall} &
\textbf{\lr{F1-score}} &
\textbf{Support} \\

\hline

$0$ & $1.00$ & $1.00$ & $1.00$ & $56651$ \\

\hline

$1$ & $0.96$ & $0.72$ & $0.82$ & $95$ \\

\hline

\textbf{Macro Avg} &
$0.98$ &
$0.86$ &
$0.91$ &
$56746$ \\

\hline

\textbf{Weighted Avg} &
$1.00$ &
$1.00$ &
$1.00$ &
$56746$ \\

\hline

\end{tabular}

\end{adjustbox}

\label{tab:knn_results}

\end{table}

ماتریس درهم‌ریختگی مدل نیز به صورت زیر به دست آمد:

\begin{table}[H]

\centering

\caption{ماتریس درهم‌ریختگی مدل KNN}

\vspace{0.15cm}

\begin{adjustbox}{max width=\columnwidth}

\begin{tabular}{|c|c|c|}

\hline

&
\textbf{پیش‌بینی کلاس 0} &
\textbf{پیش‌بینی کلاس 1} \\

\hline

\textbf{واقعی کلاس 0} &
$56648$ &
$3$ \\

\hline

\textbf{واقعی کلاس 1} &
$27$ &
$68$ \\

\hline

\end{tabular}

\end{adjustbox}

\label{tab:knn_confusion}

\end{table}

بر اساس نتایج، مدل KNN برای کلاس اقلیت دارای
\lr{Precision} برابر با $0.96$ و
\lr{Recall} برابر با $0.72$ است.

مدل توانسته است $68$ مورد از مجموع $95$ تراکنش تقلبی را
به درستی شناسایی کند، در حالی که $27$ نمونه تقلبی از دست
رفته‌اند.

همچنین تنها $3$ تراکنش سالم به اشتباه به عنوان تقلب
شناسایی شده‌اند. مقدار \lr{F1-score} برابر با $0.82$ است.

% ==================================================
% Logistic Regression
% ==================================================

\subsection{رگرسیون لجستیک}

رگرسیون لجستیک یکی از روش‌های کلاسیک برای مسائل طبقه‌بندی
دودویی است. این مدل ترکیب خطی ویژگی‌ها را محاسبه کرده و
با استفاده از تابع سیگموید احتمال تعلق نمونه به کلاس مثبت
را به دست می‌آورد.

تابع سیگموید به صورت زیر تعریف می‌شود:

\[
\sigma(z)
=
\frac{1}{1+e^{-z}}
\]

که در آن:

\[
z = w^Tx+b
\]

است.

از آنجا که رگرسیون لجستیک به مقیاس ویژگی‌ها حساس است، در
این پژوهش ویژگی‌ها با استفاده از \lr{StandardScaler}
استانداردسازی شدند.

برای انتخاب بهترین تنظیمات، از \lr{Grid Search} به همراه
\lr{Stratified K-Fold Cross Validation} استفاده شد.

دو هایپرپارامتر \lr{C} و \lr{class\_weight} مورد بررسی
قرار گرفتند.

پارامتر \lr{C} میزان معکوس قدرت منظم‌سازی مدل را مشخص می‌کند.

مقدار کوچک‌تر برای \lr{C} به معنای منظم‌سازی قوی‌تر و مقدار
بزرگ‌تر آن به معنای منظم‌سازی ضعیف‌تر است.

برای این پارامتر مقادیر
$0.01$، $0.1$، $1$، $10$ و $100$ بررسی شدند.

در پارامتر \lr{class\_weight} دو مقدار \lr{None} و
\lr{balanced} بررسی شدند.

در مجموع:

\[
5 \times 2 = 10
\]

ترکیب مختلف بررسی شد.

برای ارزیابی هر ترکیب از \lr{Stratified K-Fold Cross Validation}
با پنج بخش استفاده شد و معیار انتخاب بهترین ترکیب
\lr{F1-score} بود.

\begin{table}[H]

\centering

\caption{هایپرپارامترهای انتخاب‌شده برای رگرسیون لجستیک}

\vspace{0.15cm}

\begin{adjustbox}{max width=\columnwidth}

\begin{tabular}{|c|c|c|}

\hline

\textbf{هایپرپارامتر} &
\textbf{مقادیر بررسی‌شده} &
\textbf{مقدار انتخاب‌شده} \\

\hline

\lr{C} &
$0.01,0.1,1,10,100$ &
$10$ \\

\hline

\lr{class\_weight} &
\lr{None}, \lr{balanced} &
\lr{None} \\

\hline

\end{tabular}

\end{adjustbox}

\label{tab:lr_best_params}

\end{table}

در نهایت، ترکیب شامل \lr{C=10} و
\lr{class\_weight=None} به عنوان بهترین تنظیمات انتخاب شد.

مقدار بهترین \lr{F1-score} در فرآیند اعتبارسنجی برابر با
$0.7118$ به دست آمد.

پس از انتخاب هایپرپارامترها، مدل نهایی با استفاده از داده‌های
آموزشی استانداردشده آموزش داده شد و روی مجموعه آزمون ارزیابی
گردید.

\begin{table}[H]

\centering

\caption{نتایج رگرسیون لجستیک روی مجموعه آزمون}

\vspace{0.15cm}

\begin{adjustbox}{max width=\columnwidth}

\begin{tabular}{|c|c|c|c|c|}

\hline

\textbf{کلاس} &
\textbf{Precision} &
\textbf{Recall} &
\textbf{\lr{F1-score}} &
\textbf{Support} \\

\hline

$0$ & $1.00$ & $1.00$ & $1.00$ & $56651$ \\

\hline

$1$ & $0.85$ & $0.58$ & $0.69$ & $95$ \\

\hline

\textbf{Macro Avg} &
$0.92$ &
$0.79$ &
$0.84$ &
$56746$ \\

\hline

\textbf{Weighted Avg} &
$1.00$ &
$1.00$ &
$1.00$ &
$56746$ \\

\hline

\end{tabular}

\end{adjustbox}

\label{tab:lr_results}

\end{table}

بر اساس نتایج، رگرسیون لجستیک برای کلاس اقلیت دارای
\lr{Precision} برابر با $0.85$ و
\lr{Recall} برابر با $0.58$ است.

مقدار \lr{F1-score} برای کلاس تقلب برابر با $0.69$ به دست
آمده است. این نتیجه نشان می‌دهد که در این مجموعه‌داده،
رگرسیون لجستیک نسبت به مدل‌های غیرخطی مانند درخت تصمیم و KNN
عملکرد ضعیف‌تری در تشخیص کلاس اقلیت داشته است.

\subsection{جمع‌بندی مدل‌های پایه}

در مرحله اول، سه مدل پایه شامل درخت تصمیم، KNN و رگرسیون لجستیک
مورد بررسی قرار گرفتند. برای انتخاب هایپرپارامترهای مناسب هر مدل،
از \lr{Grid Search} همراه با \lr{Stratified K-Fold} با پنج بخش و
معیار \lr{F1-score} استفاده شد.

نتایج نشان داد که KNN با مقدار \lr{F1-score} برابر با $0.82$،
بهترین عملکرد را در میان مدل‌های پایه داشته است. این مدل با
تنظیمات $K=3$، وزن‌دهی \lr{distance} و فاصله Manhattan به
\lr{Precision} برابر با $0.96$ و \lr{Recall} برابر با $0.72$
برای کلاس تقلب دست یافت.

درخت تصمیم با \lr{Precision} برابر با $0.92$،
\lr{Recall} برابر با $0.71$ و \lr{F1-score} برابر با $0.80$
عملکرد نزدیکی به KNN داشت.

در مقابل، رگرسیون لجستیک با
\lr{Precision} برابر با $0.85$،
\lr{Recall} برابر با $0.58$ و
\lr{F1-score} برابر با $0.69$
ضعیف‌ترین عملکرد را در تشخیص کلاس تقلب نشان داد.

بنابراین، در میان مدل‌های پایه، KNN بهترین و رگرسیون لجستیک
ضعیف‌ترین عملکرد را داشته است. همچنین مشاهده می‌شود که با
وجود Accuracy بسیار بالا در هر سه مدل، معیارهای مربوط به کلاس
اقلیت تفاوت قابل توجهی دارند. بنابراین در این مسئله،
\lr{Precision}، \lr{Recall} و \lr{F1-score} معیارهای مناسب‌تری
برای مقایسه مدل‌ها هستند.

بنابراین، با توجه به محدودیت عملکرد مدل‌های پایه و ماهیت
غیرخطی و نامتوازن داده‌ها، استفاده از مدل‌های قدرتمندتر و
پیچیده‌تر مانند XGBoost و Random Forest مورد توجه قرار گرفت.

هدف از این مرحله، بررسی امکان دستیابی به عملکرد بهتر در
شناسایی کلاس اقلیت و در نهایت استفاده از این مدل‌ها در روش
\lr{Ensemble} بود.

% ==================================================
% Advanced Models
% ==================================================

\section{مدل‌های پیشرفته}

پس از بررسی مدل‌های پایه، مدل‌های قدرتمندتر مبتنی بر
درخت تصمیم بررسی شدند.

در این مرحله مدل‌های زیر مورد استفاده قرار گرفتند:

\begin{itemize}
    \item XGBoost
    \item Random Forest
    \item KNN
\end{itemize}

KNN نیز به عنوان یکی از مدل‌های مورد استفاده در Ensemble
حفظ شد تا امکان بررسی ترکیب آن با مدل‌های درختی فراهم شود.

% ==================================================
% XGBoost
% ==================================================

\subsection{XGBoost}

XGBoost یکی از الگوریتم‌های قدرتمند مبتنی بر روش
\lr{Gradient Boosting} است که از مجموعه‌ای از درخت‌های تصمیم
برای ایجاد یک مدل نهایی استفاده می‌کند.

برخلاف روش‌هایی مانند جنگل تصادفی که درخت‌ها را به صورت
مستقل آموزش می‌دهند، در XGBoost درخت‌ها به صورت ترتیبی ساخته
می‌شوند. هر درخت جدید تلاش می‌کند خطاهای مدل‌های قبلی را
کاهش دهد.

تابع کلی مدل XGBoost را می‌توان به صورت زیر نمایش داد:

\[
\hat{y}_i =
\sum_{k=1}^{K} f_k(x_i)
\]

که در آن $K$ تعداد درخت‌های موجود در مدل و $f_k$
تابع مربوط به درخت تصمیم $k$-ام است.

تابع هدف XGBoost شامل تابع خطا و جمله منظم‌سازی است:

\[
\mathcal{L} =
\sum_{i=1}^{n}
l(y_i,\hat{y}_i)
+
\sum_{k=1}^{K}
\Omega(f_k)
\]

وجود جمله منظم‌سازی به کنترل پیچیدگی مدل و کاهش احتمال
بیش‌برازش کمک می‌کند.

از آنجا که مجموعه‌داده دارای عدم توازن شدید بین کلاس سالم
و تقلبی است، از پارامتر
\lr{scale\_pos\_weight} استفاده شد.

این پارامتر با استفاده از نسبت تعداد نمونه‌های کلاس سالم به
تعداد نمونه‌های کلاس تقلب محاسبه شد:

\[
\text{scale\_pos\_weight}
=
\frac{N_0}{N_1}
\]

که در آن $N_0$ تعداد نمونه‌های کلاس سالم و $N_1$ تعداد
نمونه‌های کلاس تقلب است. بنابراین، نمونه‌های کلاس اقلیت
در فرآیند یادگیری وزن بیشتری دریافت می‌کنند.

برای انتخاب بهترین تنظیمات XGBoost از
\lr{Grid Search} به همراه
\lr{Stratified K-Fold Cross Validation} با پنج بخش استفاده شد.

معیار انتخاب بهترین مدل، \lr{F1-score} در نظر گرفته شد.

در فرآیند جستجوی هایپرپارامترها، شش پارامتر مورد بررسی قرار
گرفتند.

\begin{itemize}

    \item \lr{n\_estimators}: مقادیر $300$ و $500$

    \item \lr{max\_depth}: مقادیر $3$، $4$ و $5$

    \item \lr{learning\_rate}: مقادیر $0.03$ و $0.05$

    \item \lr{min\_child\_weight}: مقادیر $1$ و $3$

    \item \lr{subsample}: مقادیر $0.8$ و $1.0$

    \item \lr{colsample\_bytree}: مقادیر $0.8$ و $1.0$

\end{itemize}

در نتیجه تعداد کل ترکیب‌های بررسی‌شده برابر با:

\[
2 \times 3 \times 2 \times 2 \times 2 \times 2
=
96
\]

ترکیب بود.

با توجه به استفاده از پنج Fold در
\lr{Cross-Validation}، در مجموع:

\[
96 \times 5 = 480
\]

بار آموزش و ارزیابی انجام شد.

بهترین ترکیب هایپرپارامترها به صورت زیر به دست آمد:

\begin{table}[H]

\centering

\caption{هایپرپارامترهای انتخاب‌شده برای XGBoost}

\vspace{0.15cm}

\begin{adjustbox}{max width=\columnwidth}

\begin{tabular}{|c|c|}

\hline

\textbf{هایپرپارامتر} & \textbf{مقدار انتخاب‌شده} \\

\hline

\lr{n\_estimators} & $500$ \\

\hline

\lr{max\_depth} & $5$ \\

\hline

\lr{learning\_rate} & $0.05$ \\

\hline

\lr{min\_child\_weight} & $1$ \\

\hline

\lr{subsample} & $0.8$ \\

\hline

\lr{colsample\_bytree} & $1.0$ \\

\hline

\end{tabular}

\end{adjustbox}

\label{tab:xgb_best_params}

\end{table}

بهترین مقدار \lr{F1-score} در فرآیند
\lr{Cross-Validation} برابر با $0.8639$ به دست آمد.

این مقدار نشان می‌دهد که XGBoost در مرحله اعتبارسنجی توانسته
است تعادل مناسبی میان \lr{Precision} و \lr{Recall} برای کلاس
اقلیت ایجاد کند.

پس از انتخاب بهترین هایپرپارامترها، مدل نهایی با استفاده از
تمام داده‌های آموزشی آموزش داده شد و سپس روی مجموعه آزمون
ارزیابی گردید.

\begin{table}[H]

\centering

\caption{نتایج مدل XGBoost روی مجموعه آزمون}

\vspace{0.15cm}

\begin{adjustbox}{max width=\columnwidth}

\begin{tabular}{|c|c|c|c|c|}

\hline

\textbf{کلاس} &
\textbf{Precision} &
\textbf{Recall} &
\textbf{\lr{F1-score}} &
\textbf{Support} \\

\hline

$0$ & $1.00$ & $1.00$ & $1.00$ & $56651$ \\

\hline

$1$ & $0.93$ & $0.79$ & $0.85$ & $95$ \\

\hline

\textbf{Macro Avg} &
$0.96$ &
$0.89$ &
$0.93$ &
$56746$ \\

\hline

\textbf{Weighted Avg} &
$1.00$ &
$1.00$ &
$1.00$ &
$56746$ \\

\hline

\end{tabular}

\end{adjustbox}

\label{tab:xgb_results}

\end{table}

بر اساس نتایج، XGBoost برای کلاس اقلیت دارای
\lr{Precision} برابر با $0.93$،
\lr{Recall} برابر با $0.79$ و
\lr{F1-score} برابر با $0.85$ است.

بنابراین مدل توانسته است حدود $79$ درصد از تراکنش‌های تقلبی
موجود در مجموعه آزمون را شناسایی کند.

همچنین مقدار \lr{Precision} برابر با $0.93$ نشان می‌دهد که
بخش عمده تراکنش‌هایی که توسط مدل به عنوان تقلب شناسایی شده‌اند،
واقعاً تقلبی بوده‌اند.

مقدار \lr{F1-score} برابر با $0.85$ نیز نشان‌دهنده تعادل
مناسب میان \lr{Precision} و \lr{Recall} در تشخیص کلاس تقلب است.

با توجه به اینکه معیار انتخاب مدل در فرآیند
\lr{Grid Search}، مقدار \lr{F1-score} بوده است، نتیجه
اعتبارسنجی $0.8639$ و نتیجه آزمون $0.85$ به دست آمده است.

نزدیک بودن این دو مقدار نشان می‌دهد که افت عملکرد مدل از
مرحله اعتبارسنجی به مجموعه آزمون نسبتاً محدود بوده است.

در مقایسه با مدل‌های پایه، XGBoost عملکرد بهتری در تشخیص
کلاس اقلیت نشان داده است. به طور مشخص، مقدار
\lr{F1-score} این مدل برای کلاس تقلب $0.85$ است که نسبت به
درخت تصمیم با مقدار $0.80$ و رگرسیون لجستیک با مقدار
$0.69$ بیشتر است و با KNN با مقدار $0.82$ نیز عملکرد
بهتری دارد.

همچنین مقدار \lr{Precision} برابر با $0.93$ در کنار
\lr{Recall} برابر با $0.79$ نشان می‌دهد که XGBoost گزینه
مناسبی برای استفاده در مرحله بعدی، یعنی ترکیب مدل‌ها در
روش \lr{Ensemble} است.

% ==================================================
% Random Forest
% ==================================================

\subsection{جنگل تصادفی}

جنگل تصادفی از تعداد زیادی درخت تصمیم تشکیل شده است.
هر درخت روی نمونه‌ها و ویژگی‌های متفاوتی آموزش داده می‌شود
و در نهایت خروجی درخت‌ها با یکدیگر ترکیب می‌شوند.

احتمال تعلق یک نمونه به کلاس تقلب را می‌توان به صورت زیر
در نظر گرفت:

\[
P(y=1|x)
=
\frac{1}{T}
\sum_{t=1}^{T}
P_t(y=1|x)
\]

که در آن $T$ تعداد درخت‌ها است.

مزیت مهم جنگل تصادفی کاهش واریانس و کاهش احتمال
بیش‌برازش نسبت به یک درخت تصمیم منفرد است.

در این پژوهش پارامترهایی مانند تعداد درخت‌ها،
\lr{max\_depth}، \lr{min\_samples\_split}،
\lr{min\_samples\_leaf} و \lr{max\_features} مورد بررسی
قرار گرفتند.

برای انتخاب تنظیمات مناسب، از Stratified K-Fold
Cross-Validation و معیار \lr{F1-score} استفاده شد.

با توجه به عدم توازن شدید بین کلاس سالم و کلاس تقلب، در این
پژوهش از پارامتر \lr{class\_weight=balanced} استفاده شد.

برای انتخاب بهترین تنظیمات جنگل تصادفی از
\lr{Grid Search} همراه با
\lr{Stratified K-Fold Cross Validation} با پنج بخش استفاده شد.

معیار انتخاب بهترین مدل، \lr{F1-score} در نظر گرفته شد.

در فرآیند جستجوی هایپرپارامترها، پنج پارامتر مورد بررسی قرار
گرفتند.

\begin{itemize}

\item \lr{n\_estimators}: مقادیر $200$ و $400$

\item \lr{max\_depth}: مقادیر $10$، $11$، $12$، $13$ و $14$

\item \lr{min\_samples\_split}: مقادیر $2$ و $5$

\item \lr{min\_samples\_leaf}: مقادیر $1$ و $2$

\item \lr{max\_features}: مقدار \lr{sqrt}

\end{itemize}

بنابراین تعداد کل ترکیب‌های بررسی‌شده برابر با:

\[
2 \times 5 \times 2 \times 2 \times 1 = 40
\]

ترکیب بود.

با توجه به استفاده از پنج Fold در
\lr{Cross-Validation}، در مجموع:

\[
40 \times 5 = 200
\]

بار آموزش و ارزیابی انجام شد.

بهترین ترکیب هایپرپارامترها به صورت زیر به دست آمد:

\begin{table}[H]

\centering

\caption{هایپرپارامترهای انتخاب‌شده برای جنگل تصادفی}

\vspace{0.15cm}

\begin{adjustbox}{max width=\columnwidth}

\begin{tabular}{|c|c|}

\hline

\textbf{هایپرپارامتر} & \textbf{مقدار انتخاب‌شده} \\

\hline

\lr{n\_estimators} & $200$ \\

\hline

\lr{max\_depth} & $13$ \\

\hline

\lr{min\_samples\_split} & $2$ \\

\hline

\lr{min\_samples\_leaf} & $2$ \\

\hline

\lr{max\_features} & \lr{sqrt} \\

\hline

\end{tabular}

\end{adjustbox}

\label{tab:rf_best_params}

\end{table}

بهترین مقدار \lr{F1-score} در فرآیند
\lr{Cross-Validation} برابر با $0.8495$ به دست آمد.

پس از انتخاب بهترین هایپرپارامترها، مدل نهایی با استفاده از
تمام داده‌های آموزشی آموزش داده شد و سپس روی مجموعه آزمون
ارزیابی گردید.

\begin{table}[H]

\centering

\caption{نتایج مدل جنگل تصادفی روی مجموعه آزمون}

\vspace{0.15cm}

\begin{adjustbox}{max width=\columnwidth}

\begin{tabular}{|c|c|c|c|c|}

\hline

\textbf{کلاس} &
\textbf{Precision} &
\textbf{Recall} &
\textbf{\lr{F1-score}} &
\textbf{Support} \\

\hline

$0$ & $1.00$ & $1.00$ & $1.00$ & $56651$ \\

\hline

$1$ & $0.92$ & $0.73$ & $0.81$ & $95$ \\

\hline

\textbf{Macro Avg} &
$0.96$ &
$0.86$ &
$0.91$ &
$56746$ \\

\hline

\textbf{Weighted Avg} &
$1.00$ &
$1.00$ &
$1.00$ &
$56746$ \\

\hline

\end{tabular}

\end{adjustbox}

\label{tab:rf_results}

\end{table}

بر اساس نتایج، جنگل تصادفی برای کلاس اقلیت دارای
\lr{Precision} برابر با $0.92$،
\lr{Recall} برابر با $0.73$ و
\lr{F1-score} برابر با $0.81$ است.

مدل توانسته است حدود $73$ درصد از تراکنش‌های تقلبی موجود
در مجموعه آزمون را شناسایی کند.

مقدار \lr{Precision} برابر با $0.92$ نشان می‌دهد که مدل در
تشخیص تراکنش‌های تقلبی دقت مناسبی دارد.

مقدار \lr{F1-score} برابر با $0.81$ نیز نشان‌دهنده تعادل
مناسب میان \lr{Precision} و \lr{Recall} برای کلاس اقلیت است.

مقایسه مقدار \lr{F1-score} حاصل از Cross-Validation با نتیجه
مجموعه آزمون نشان می‌دهد که مقدار این معیار از $0.8495$ در
فرآیند اعتبارسنجی به $0.81$ در مجموعه آزمون کاهش یافته است.

این اختلاف نشان می‌دهد که عملکرد مدل روی داده‌های آزمون
کمی پایین‌تر از عملکرد متوسط آن در فرآیند اعتبارسنجی بوده است.

در مقایسه با XGBoost، جنگل تصادفی در این آزمایش عملکرد
ضعیف‌تری در تشخیص کلاس تقلب داشته است. XGBoost با
\lr{Precision} برابر با $0.93$،
\lr{Recall} برابر با $0.79$ و
\lr{F1-score} برابر با $0.85$ عملکرد بهتری نسبت به جنگل
تصادفی با مقادیر $0.92$، $0.73$ و $0.81$ داشته است.

با این حال، جنگل تصادفی همچنان عملکرد مناسبی دارد و به دلیل
تفاوت در ساختار یادگیری آن با XGBoost، می‌تواند به عنوان یکی
از مدل‌های مناسب برای مرحله بعدی یعنی روش
\lr{Ensemble} مورد استفاده قرار گیرد.

% ==================================================
% Ensemble OOF
% ==================================================

\subsection{انتخاب وزن‌های Ensemble با استفاده از OOF}

پس از آموزش مدل‌های XGBoost، Random Forest و KNN، در این مرحله
هدف انتخاب وزن مناسب برای ترکیب این سه مدل بود.

برای جلوگیری از انتخاب وزن‌ها بر اساس پیش‌بینی‌هایی که مدل در
فرآیند آموزش خود مشاهده کرده است، از روش
\lr{Out-of-Fold Predictions} یا به اختصار \lr{OOF} استفاده شد.

در این روش از \lr{Stratified K-Fold Cross Validation} با پنج
بخش استفاده شد.

در هر Fold، مدل‌ها روی چهار بخش آموزش داده شده و روی بخش
باقی‌مانده پیش‌بینی انجام دادند. بنابراین برای هر نمونه آموزشی،
یک احتمال پیش‌بینی‌شده توسط مدلی به دست آمد که آن نمونه را
در مرحله آموزش مشاهده نکرده بود.

برای هر سه مدل، احتمال تعلق نمونه به کلاس تقلب ذخیره شد:

\[
P_{\mathrm{XGB}}^{\mathrm{OOF}},
\qquad
P_{\mathrm{RF}}^{\mathrm{OOF}},
\qquad
P_{\mathrm{KNN}}^{\mathrm{OOF}}
\]

سپس احتمال نهایی Ensemble به صورت ترکیب وزنی احتمال‌های
سه مدل محاسبه شد:

\[
P_{\mathrm{ensemble}}
=
w_{\mathrm{XGB}}P_{\mathrm{XGB}}
+
w_{\mathrm{RF}}P_{\mathrm{RF}}
+
w_{\mathrm{KNN}}P_{\mathrm{KNN}}
\]

که در آن مجموع وزن‌ها برابر با یک است:

\[
w_{\mathrm{XGB}}
+
w_{\mathrm{RF}}
+
w_{\mathrm{KNN}}
=
1
\]

برای انتخاب وزن‌ها، مقدار وزن XGBoost و Random Forest با
گام $0.05$ بررسی شد و وزن KNN به صورت زیر محاسبه گردید:

\[
w_{\mathrm{KNN}}
=
1-w_{\mathrm{XGB}}-w_{\mathrm{RF}}
\]

همچنین برای احتمال نهایی، آستانه‌های مختلف از $0.05$ تا
$0.95$ با گام $0.01$ بررسی شدند.

برای هر ترکیب وزن و آستانه، معیارهای
\lr{Precision}، \lr{Recall} و \lr{F1-score} محاسبه شدند.

هدف اصلی انتخاب ترکیبی بود که شرایط زیر را برآورده کند:

\[
\mathrm{Precision} \geq 0.90
\]

و

\[
\mathrm{Recall} \geq 0.80
\]

در میان ترکیب‌هایی که این دو شرط را برآورده می‌کردند،
ترکیبی انتخاب شد که بیشترین \lr{F1-score} را داشته باشد.

در صورتی که هیچ ترکیبی نتواند به صورت همزمان دو شرط فوق را
برآورده کند، ترکیبی که بیشترین \lr{F1-score} را در میان تمام
ترکیب‌های بررسی‌شده داشته باشد، به عنوان ترکیب نهایی انتخاب
می‌شود.

روند انتخاب وزن‌ها و آستانه به صورت الگوریتمی انجام شد و
پیش‌بینی‌های \lr{OOF} هر سه مدل برای تمام نمونه‌های آموزشی
تولید شدند.

پس از اتمام فرآیند جستجو، وزن‌های نهایی سه مدل و آستانه
تصمیم به صورت زیر گزارش شدند:

\begin{table}[H]

\centering

\caption{پارامترهای انتخاب‌شده برای Ensemble}

\vspace{0.15cm}

\begin{adjustbox}{max width=\columnwidth}

\begin{tabular}{|c|c|}

\hline

\textbf{پارامتر} & \textbf{مقدار} \\

\hline

وزن XGBoost & $0.45$ \\

\hline

وزن Random Forest & $0.45$ \\

\hline

وزن KNN & $0.10$ \\

\hline

آستانه تصمیم & $0.56$ \\

\hline

\lr{F1-score} روی OOF & $0.8727$ \\

\hline

\end{tabular}

\end{adjustbox}

\label{tab:ensemble_oof_params}

\end{table}

مقادیر نهایی جدول فوق مستقیماً از اجرای فرآیند
\lr{OOF} و جستجوی وزن‌ها و آستانه استخراج می‌شوند.

نکته مهم این است که در این مرحله، داده‌های مجموعه آزمون
برای انتخاب وزن‌ها و آستانه استفاده نشده‌اند. بنابراین
مجموعه آزمون همچنان به عنوان یک مجموعه مستقل برای ارزیابی
نهایی عملکرد Ensemble حفظ می‌شود.

% ==================================================
% Ensemble Results
% ==================================================

\subsection{نتایج Ensemble}

در مرحله نهایی، احتمال‌های تولیدشده توسط سه مدل XGBoost،
Random Forest و KNN با استفاده از یک ترکیب وزنی با یکدیگر
ادغام شدند.

برای انتخاب وزن‌های مناسب، از
\lr{Out-of-Fold Predictions} حاصل از
\lr{Stratified K-Fold Cross Validation} با پنج بخش استفاده شد.

در فرآیند جستجو، وزن‌های XGBoost و Random Forest با گام
$0.05$ بررسی شدند و وزن KNN به گونه‌ای تعیین شد که مجموع
وزن‌های سه مدل برابر با یک باشد.

همچنین آستانه‌های مختلف بین $0.05$ و $0.95$ با گام $0.01$
مورد بررسی قرار گرفتند.

هدف، انتخاب ترکیبی بود که در صورت امکان دارای
\lr{Precision} حداقل $0.90$ و
\lr{Recall} حداقل $0.80$ باشد و در میان ترکیب‌های مجاز،
بیشترین \lr{F1-score} را ایجاد کند.

در نهایت، بهترین ترکیب بر اساس پیش‌بینی‌های
\lr{OOF} به صورت زیر به دست آمد:

\[
w_{\mathrm{XGB}} = 0.45
\]

\[
w_{\mathrm{RF}} = 0.45
\]

\[
w_{\mathrm{KNN}} = 0.10
\]

همچنین آستانه تصمیم برابر با:

\[
\tau = 0.56
\]

انتخاب شد.

بنابراین احتمال نهایی Ensemble به صورت زیر محاسبه می‌شود:

\[
P_{\mathrm{ensemble}}
=
0.45P_{\mathrm{XGB}}
+
0.45P_{\mathrm{RF}}
+
0.10P_{\mathrm{KNN}}
\]

و بر اساس آستانه $0.56$، برچسب نهایی نمونه‌ها تعیین می‌شود:

\[
\hat{y}
=
\begin{cases}
1 & P_{\mathrm{ensemble}}\geq0.56\\
0 & P_{\mathrm{ensemble}}<0.56
\end{cases}
\]

مقدار \lr{F1-score} به دست آمده از پیش‌بینی‌های
\lr{OOF} برابر با $0.8727$ بود.

نتایج نهایی مدل Ensemble روی مجموعه آزمون به صورت زیر
به دست آمد:

\begin{table}[H]

\centering

\caption{نتایج نهایی Ensemble روی مجموعه آزمون}

\vspace{0.15cm}

\begin{adjustbox}{max width=\columnwidth}

\begin{tabular}{|c|c|}

\hline

\textbf{پارامتر} & \textbf{مقدار} \\

\hline

وزن XGBoost & $0.45$ \\

\hline

وزن Random Forest & $0.45$ \\

\hline

وزن KNN & $0.10$ \\

\hline

Threshold & $0.56$ \\

\hline

Precision & $0.9730$ \\

\hline

Recall & $0.7579$ \\

\hline

\lr{F1-score} & $0.8521$ \\

\hline

PR-AUC & $0.8253$ \\

\hline

\end{tabular}

\end{adjustbox}

\label{tab:ensemble-results}

\end{table}

همان‌طور که مشاهده می‌شود، مدل Ensemble روی مجموعه آزمون
به \lr{Precision} برابر با $0.9730$ دست یافته است.

مقدار \lr{Recall} برابر با $0.7579$ است. بنابراین مدل توانسته
است $72$ مورد از مجموع $95$ تراکنش تقلبی موجود در مجموعه آزمون
را به درستی شناسایی کند.

مقدار \lr{F1-score} نیز برابر با $0.8521$ به دست آمده است که
نشان‌دهنده تعادل مناسبی میان \lr{Precision} و \lr{Recall} است.

ماتریس درهم‌ریختگی مدل Ensemble به صورت زیر به دست آمد:

\begin{table}[H]

\centering

\caption{ماتریس درهم‌ریختگی Ensemble}

\vspace{0.15cm}

\begin{adjustbox}{max width=\columnwidth}

\begin{tabular}{|c|c|c|}

\hline

&
\textbf{پیش‌بینی کلاس 0} &
\textbf{پیش‌بینی کلاس 1} \\

\hline

\textbf{واقعی کلاس 0} &
$56649$ &
$2$ \\

\hline

\textbf{واقعی کلاس 1} &
$23$ &
$72$ \\

\hline

\end{tabular}

\end{adjustbox}

\label{tab:ensemble-confusion}

\end{table}

بر اساس ماتریس درهم‌ریختگی، مدل تنها $2$ نمونه سالم را به
اشتباه به عنوان تقلب شناسایی کرده است. در مقابل، از مجموع
$95$ نمونه تقلبی، $72$ نمونه به درستی شناسایی شده و $23$
نمونه از دست رفته‌اند.

مقدار \lr{PR-AUC} مدل نیز برابر با $0.8253$ به دست آمد.

این معیار با توجه به نامتوازن بودن شدید داده‌ها، معیار
مناسبی برای ارزیابی عملکرد مدل در تشخیص کلاس اقلیت محسوب
می‌شود.

در مقایسه با مدل‌های منفرد، Ensemble توانسته است
\lr{Precision} بسیار بالایی ایجاد کند. برای مثال، XGBoost
به تنهایی دارای \lr{Precision} برابر با $0.93$،
\lr{Recall} برابر با $0.79$ و
\lr{F1-score} برابر با $0.85$ بود، در حالی که Ensemble به
\lr{Precision} برابر با $0.97$ و
\lr{F1-score} برابر با $0.85$ دست یافت.

با این حال، \lr{Recall} Ensemble برابر با $0.76$ است که
نسبت به Recall برابر با $0.79$ در XGBoost کمتر است.

بنابراین ترکیب مدل‌ها در این آزمایش باعث افزایش Precision
و کاهش تعداد مثبت‌های کاذب شده است، اما به قیمت کاهش نسبی
Recall تمام شده است.

همچنین مقدار \lr{F1-score} حاصل از OOF برابر با $0.8727$ و
مقدار آن روی مجموعه آزمون برابر با $0.8521$ است.

کاهش محدود این معیار از مرحله OOF به مجموعه آزمون نشان می‌دهد
که عملکرد مدل روی داده‌های دیده‌نشده همچنان در سطح مناسبی
باقی مانده است.

با توجه به اینکه هدف اولیه انتخاب ترکیب شامل
\lr{Precision} حداقل $0.90$ و
\lr{Recall} حداقل $0.80$ بود، ترکیب انتخاب‌شده روی مجموعه
آزمون نتوانسته است شرط Recall برابر با $0.80$ را حفظ کند.

بنابراین می‌توان گفت که این ترکیب در مجموعه آموزش و
پیش‌بینی‌های OOF عملکرد بسیار خوبی داشته، اما روی مجموعه
آزمون Recall آن به حدود $0.76$ کاهش یافته است.

% ==================================================
% Overall Results
% ==================================================

\section{نتایج کلی و تحلیل مقایسه‌ای}

در این بخش، عملکرد تمام مدل‌های بررسی‌شده با یکدیگر مقایسه
می‌شود. با توجه به نامتوازن بودن شدید مجموعه‌داده، تمرکز اصلی
بر روی عملکرد مدل در شناسایی کلاس اقلیت، یعنی تراکنش‌های تقلبی،
قرار گرفته است.

بنابراین معیارهای \lr{Precision}،
\lr{Recall} و \lr{F1-score} نسبت به \lr{Accuracy} اهمیت
بیشتری دارند.

نتایج نهایی مدل‌ها در جدول زیر خلاصه شده است:

\begin{table}[H]

\centering

\caption{مقایسه عملکرد مدل‌های مختلف روی مجموعه آزمون}

\vspace{0.15cm}

\begin{adjustbox}{max width=\columnwidth}

\begin{tabular}{|l|c|c|c|}

\hline

\textbf{مدل} &
\textbf{Precision} &
\textbf{Recall} &
\textbf{\lr{F1-score}} \\

\hline

Decision Tree & $0.92$ & $0.71$ & $0.80$ \\

\hline

KNN & $0.96$ & $0.72$ & $0.82$ \\

\hline

Logistic Regression & $0.85$ & $0.58$ & $0.69$ \\

\hline

Random Forest & $0.92$ & $0.73$ & $0.81$ \\

\hline

XGBoost & $0.93$ & $0.79$ & $0.85$ \\

\hline

Ensemble & $0.97$ & $0.76$ & $0.85$ \\

\hline

\end{tabular}

\end{adjustbox}

\label{tab:overall-results}

\end{table}

با توجه به نتایج جدول، رگرسیون لجستیک با
\lr{F1-score} برابر با $0.69$ ضعیف‌ترین عملکرد را در میان
مدل‌های بررسی‌شده داشته است.

همچنین \lr{Recall} برابر با $0.58$ نشان می‌دهد که این مدل
بخش قابل توجهی از تراکنش‌های تقلبی را از دست داده است.

این نتیجه نشان می‌دهد که مدل‌های خطی در این آزمایش توانایی
کمتری برای مدل‌سازی روابط پیچیده و غیرخطی موجود در داده‌ها
داشته‌اند.

در میان مدل‌های پایه، KNN بهترین عملکرد را با
\lr{F1-score} برابر با $0.82$ نشان داده است.

این مدل همچنین به \lr{Precision} برابر با $0.96$ دست یافته است؛
بنابراین تعداد نمونه‌های سالمی که به اشتباه به عنوان تقلب
شناسایی شده‌اند، نسبتاً کم بوده است.

جنگل تصادفی با \lr{F1-score} برابر با $0.81$ عملکردی نزدیک
به KNN داشته است.

با این حال، \lr{Recall} برابر با $0.73$ نشان می‌دهد که این
مدل نسبت به XGBoost تعداد بیشتری از تراکنش‌های تقلبی را
از دست داده است.

بهترین عملکرد در میان مدل‌های منفرد مربوط به XGBoost بوده
است. این مدل با \lr{Precision} برابر با $0.93$،
\lr{Recall} برابر با $0.79$ و
\lr{F1-score} برابر با $0.85$ توانسته است بهترین تعادل
میان تشخیص صحیح تراکنش‌های تقلبی و جلوگیری از ایجاد
مثبت‌های کاذب را ایجاد کند.

مقایسه XGBoost با Random Forest نیز نشان می‌دهد که XGBoost
در هر سه معیار اصلی عملکرد بهتری داشته است.

مقدار \lr{Recall} از $0.73$ در Random Forest به $0.79$ در
XGBoost افزایش یافته و \lr{F1-score} نیز از $0.81$ به $0.85$
رسیده است.

این نتیجه نشان می‌دهد که ساختار ترتیبی و فرآیند Boosting در
XGBoost در این مجموعه‌داده توانسته است روابط پیچیده‌تری را
نسبت به روش Bagging جنگل تصادفی مدل کند.

در مرحله نهایی، سه مدل XGBoost، Random Forest و KNN با
استفاده از یک Ensemble وزنی ترکیب شدند.

وزن‌های انتخاب‌شده شامل $0.45$ برای XGBoost،
$0.45$ برای Random Forest و $0.10$ برای KNN بودند.

همچنین آستانه تصمیم برابر با $0.56$ انتخاب شد.

Ensemble توانست \lr{Precision} را از $0.93$ در XGBoost
به $0.97$ افزایش دهد.

همچنین تعداد مثبت‌های کاذب در مجموعه آزمون تنها $2$ مورد بود.

این موضوع نشان می‌دهد که ترکیب مدل‌ها در کاهش
\lr{False Positive} عملکرد بسیار خوبی داشته است.

با این حال، افزایش Precision با کاهش Recall همراه بوده است.

Recall مدل Ensemble برابر با $0.76$ بوده، در حالی که این
مقدار برای XGBoost برابر با $0.79$ است.

بنابراین Ensemble در آزمایش حاضر نتوانسته است از نظر
شناسایی تعداد بیشتری از تراکنش‌های تقلبی نسبت به XGBoost
عملکرد بهتری داشته باشد.

مقدار \lr{F1-score} هر دو مدل XGBoost و Ensemble برابر با
تقریباً $0.85$ است.

بنابراین بر اساس معیار F1، ترکیب مدل‌ها مزیت مشخصی نسبت به
XGBoost به تنهایی ایجاد نکرده است.

با این حال، Ensemble زمانی مزیت بیشتری دارد که کاهش
مثبت‌های کاذب و افزایش Precision اهمیت بیشتری داشته باشد.

یکی دیگر از نتایج مهم این آزمایش، تفاوت میان عملکرد OOF و
عملکرد مجموعه آزمون است.

مقدار \lr{F1-score} برای Ensemble در پیش‌بینی‌های OOF برابر
با $0.8727$ بود، در حالی که این مقدار روی مجموعه آزمون به
$0.8521$ رسید.

این کاهش نشان می‌دهد که عملکرد مدل روی داده‌های کاملاً
دیده‌نشده کمی پایین‌تر از عملکرد آن در فرآیند انتخاب وزن‌ها
بوده است.

در مجموع، نتایج نشان می‌دهد که انتخاب بهترین مدل به هدف
سیستم تشخیص تقلب بستگی دارد.

اگر هدف اصلی ایجاد تعادل میان \lr{Precision} و \lr{Recall}
و شناسایی تعداد بیشتری از تراکنش‌های تقلبی باشد، XGBoost
در این آزمایش بهترین گزینه است.

از طرف دیگر، اگر کاهش مثبت‌های کاذب اهمیت بیشتری داشته باشد،
Ensemble با \lr{Precision} برابر با $0.97$ گزینه مناسب‌تری
خواهد بود.

بنابراین، بر اساس نتایج مجموعه آزمون، XGBoost را می‌توان
به عنوان بهترین مدل منفرد این پژوهش در نظر گرفت، در حالی
که Ensemble نشان داده است که ترکیب مدل‌ها لزوماً بهبود
همزمان تمام معیارهای ارزیابی را تضمین نمی‌کند و انتخاب
وزن‌ها و آستانه تصمیم نقش مهمی در تعیین نقطه تعادل میان
\lr{Precision} و \lr{Recall} دارد.



\begin{thebibliography}{9}

\bibitem{han}
J. Han, M. Kamber, and J. Pei,
\textit{Data Mining: Concepts and Techniques}.

\bibitem{goodfellow}
I. Goodfellow, Y. Bengio, and A. Courville,
\textit{Deep Learning},
MIT Press, 2016.

\bibitem{bishop}
C. M. Bishop,
\textit{Pattern Recognition and Machine Learning},
Springer, 2006.

\end{thebibliography}

\end{document}